\documentclass[11pt]{article}
\usepackage[margin=1in]{geometry}
\usepackage{amsmath,amssymb,amsthm,hyperref}
\newtheorem{theorem}{Theorem}
\newtheorem{lemma}{Lemma}
\title{Strict log-concavity of the gaps between squared Bessel zeros\\
\large A solution to MathDB 326635}
\author{Principia Math in Progress}
\date{12 September 2026}
\begin{document}
\maketitle
\begin{abstract}
The ratios of consecutive gaps between the squared positive zeros of $J_1$,
with the additional zero $j_{1,0}=0$, are strictly decreasing. This proves
Conjecture 3 of Brix, Canuto, and Dahmen. The proof uses an affine comparison
for the equation $u''+u/(4z)=0$ and includes the singular endpoint at zero.
\end{abstract}

\section*{Task 1: solve the original problem}
Let $j_{1,k}$ be the $k$th positive zero of $J_1$, set $j_{1,0}=0$, and put
\[
 \lambda_k=j_{1,k}^2,\qquad
 d_k=\lambda_k-\lambda_{k-1},\qquad
 q_k=\frac{d_{k+1}}{d_k}\quad(k\geq1).
\]
Brix--Canuto--Dahmen \cite{BCD}, Conjecture 3, ask whether $q_k$ decreases.
The same question is explicitly retained as Conjecture 3.36, equation
(3.51), in Brix's later thesis \cite{Brix}, page 38.

\begin{theorem}\label{main}
For every $k\geq1$,
\[
 d_{k+1}>d_k>0,\qquad
 d_{k+1}^2>d_kd_{k+2}.
\]
Equivalently,
\[
 q_1>q_2>q_3>\cdots>1.
\]
Moreover, $q_k\to1$, and the sharp upper bound on the sequence is
$q_1=j_{1,2}^2/j_{1,1}^2-1$.
\end{theorem}

\subsection*{The equation for the squared zeros}
The entire function
\[
 u(z)=2\sqrt z\,J_1(\sqrt z)
 =z\sum_{n=0}^{\infty}\frac{(-z/4)^n}{n!(n+1)!}
\]
is real on the real axis and satisfies
\begin{equation}\label{ode}
 u''(z)+\frac{u(z)}{4z}=0\qquad(z>0).
\end{equation}
This follows directly from its power series or the Bessel equation.
Its nonnegative zeros are exactly $\lambda_0,\lambda_1,\ldots$,
all simple; in particular
\[
 u(z)=z+O(z^2),\qquad u'(0)=1.
\]
Simplicity at a positive zero also follows from uniqueness for \eqref{ode}.

\subsection*{Comparing two successive intervals}
Fix $k\geq1$, and write
\[
 a=\lambda_{k-1},\quad b=\lambda_k,\quad c=\lambda_{k+1},
 \quad d=\lambda_{k+2},\qquad
 r=\frac{c-b}{b-a},\quad s=b-ra.
\]
Thus $r>0$, and the affine function $T(z)=rz+s$ sends $a$ to $b$
and $b$ to $c$. Multiply $u$ by a sign, if necessary, so that it is
positive on $(a,b)$. Since its signs alternate at consecutive zeros,
\[
 v(z)=-u(T(z))
\]
is also positive on $(a,b)$, and both $u$ and $v$ vanish at its endpoints.
The function $v$ solves
\[
 v''(z)+\frac{r^2}{4(rz+s)}v(z)=0.
\]
Set
\begin{equation}\label{D}
 D(z)=\frac{r^2}{4(rz+s)}-\frac1{4z}
 =\frac{r(r-1)z-s}{4z(rz+s)}.
\end{equation}
For $z>a$ the denominator is positive. With
$W(z)=u'(z)v(z)-u(z)v'(z)$, differentiation gives
\begin{equation}\label{wronskian}
 W'(z)=D(z)u(z)v(z).
\end{equation}
Both endpoint Wronskians vanish, hence
\begin{equation}\label{balance}
 \int_a^b D(z)u(z)v(z)\,dz=0.
\end{equation}
When $a=0$, this is an improper integral justified without any regular
endpoint assumption: $u(z)=O(z)$, $v(z)=O(z)$, their derivatives are
bounded, and $D(z)=O(1/z)$. Thus the integrand is $O(z)$ and
$W(z)\to0$ as $z\downarrow0$.

If $r\leq1$, then $s=b-ra\geq b-a>0$, so the numerator in \eqref{D}
is strictly negative throughout $(a,b)$. This contradicts \eqref{balance},
because $uv>0$ there. Therefore $r>1$.
The numerator of \eqref{D} is now a strictly increasing affine function.
Its integral against the positive weight $uv/[4z(rz+s)]$ is zero, so
it must change sign at a point strictly between $a$ and $b$.
In particular,
\begin{equation}\label{positive}
 D(z)>0\qquad(z\geq b).
\end{equation}

The next zero of $v$ after $b$ is $e=T^{-1}(d)$.
We claim $e<c$. Otherwise $u$ and $v$ are both negative on $(b,c)$.
Equations \eqref{wronskian} and \eqref{positive} give $W(c)>W(b)=0$.
But $u(c)=0$, $u'(c)>0$, and $v(c)\leq0$, so
\[
 W(c)=u'(c)v(c)\leq0,
\]
a contradiction. Hence $e<c$, or $d<T(c)$. Since $T(b)=c$, this is
\[
 d-c<r(c-b).
\]
It follows that
\[
 1<r=\frac{c-b}{b-a},\qquad
 \frac{d-c}{c-b}<\frac{c-b}{b-a},
\]
which proves both strict inequalities in Theorem \ref{main}, for every $k$.

\subsection*{The limit and the sharp bound}
The decreasing sequence $q_k>1$ has a limit $L\geq1$.
If $L>1$, then $d_{k+1}\geq Ld_k$, so the gaps, and consequently
$\lambda_k$, grow at least geometrically.
This contradicts the classical Bessel zero asymptotic
$j_{1,k}/k\to\pi$ \cite{DLMF}. Thus $L=1$.
The strict decrease makes the first term the attained maximum,
proving the sharp upper bound as well.

\section*{Follow-up tasks after the full solution}
\noindent\textbf{Task 2.}
Obtain explicit bounds for $q_k-1$ with optimal leading constants and
rigorous remainders uniform in $k$.

\medskip
\noindent\textbf{Task 3.}
Determine which exponents $p$ make the ratios of consecutive zero
spacings of $y''+\kappa z^{-p}y=0$ decrease, remain constant, or increase,
including the relevant singular endpoint conditions.

\medskip
\noindent\textbf{Task 4.}
Classify the real orders $\nu$ for which the positive gaps
$j_{\nu,k}^2-j_{\nu,k-1}^2$ are log-concave, specifying separately any
adjoined zero at $k=0$.

\medskip
\noindent\textbf{Task 5.}
Use the integral balance \eqref{balance} to test whether the analogous
affine comparison closes the LGL interval-ratio conjecture in the same
source, including the dependence on polynomial degree.

\section*{Verification and status}
The accompanying verifier checks the differential equation, affine
transformation, potential difference, Wronskian identity, and gap-ratio
algebra symbolically. Supplementary zero calculations test the strict
inequalities and the location of the potential crossing; they are not
the proof of the infinite assertions.

MathDB was checked live before the attempt and again before registration:
the record was open and its claim fields were empty. The original source,
the author's 2018 thesis explicitly retaining the exact conjecture, both
indexed citing works, and public claim catalogues were checked.
The literature audit records the sources, dates, and access limits.

\begin{thebibliography}{9}
\bibitem{BCD}
K. Brix, C. Canuto, and W. Dahmen,
\emph{Nested dyadic grids associated with Legendre--Gauss--Lobatto grids},
Numerische Mathematik \textbf{131} (2015), 205--239,
doi:10.1007/s00211-014-0691-4.
Author preprint: \emph{Legendre-Gauss-Lobatto grids and associated nested
dyadic grids}, arXiv:1311.0028v1 (2013).

\bibitem{Brix}
K. Brix, \emph{Robust Preconditioners for hp-Discontinuous Galerkin
Discretizations for Elliptic Problems}, doctoral thesis, RWTH Aachen,
2018; repository publication 2019, doi:10.18154/RWTH-2019-10662.
Conjecture 3.36, p.~38; its reference [34] is the final journal paper above.
\url{https://publications.rwth-aachen.de/record/772363}.

\bibitem{DLMF}
NIST Digital Library of Mathematical Functions,
\S10.21(vi), equation (10.21.19),
\url{https://dlmf.nist.gov/10.21#vi}.
\end{thebibliography}
\end{document}
